%% title: The trap split: adaptation set, probe set, held-out species
%% target: Sections/Methods.tex
%% anchor: \subsubsection{Light-trap images}\label{sec:data_traps}
%% action: after-paragraph
%% why: The same trap images are both the benchmark and the only source of unlabelled target images. Without this split, self-training would train on its own test set. Every trap result in the paper uses it.
%% source: data/flemming/splits/split_summary.json; journal/2026-08-02-the-shifted-benchmark-is-also-the-adaptation-set.md; journal/2026-10-12-do-trap-moths-leak-across-nights.md
%% check: The 0.0003 was measured in the old metric convention; re-check after rescoring (TODO in the comment). | Asger's calibration split is random over crops; this paragraph is the argument for grouping it too.
%% status: ready
Because the trap images serve both as the benchmark and as the only source of unlabelled target-domain images (\cref{sec:selftraining}), they are partitioned by capture group, a pair of trap and night with nights running from midday to midday.
Groups are assigned at random to an adaptation set (27,230 crops), a probe set (15,200 crops of 368 species) that no training run touches, and an unused remainder (5,475 crops).
Fifteen percent of the species (75) are withheld from adaptation altogether.
Their probe crops (2,455 crops of 58 species) form the held-out-species subset, which separates adaptation to the capture conditions from specialisation on the species that were adapted to.
Grouping by night is necessary: 35\,\% of crops have a near pixel-identical twin elsewhere in the same night.
It is also nearly sufficient: 13 probe crops (0.09\,\%) have a twin in the adaptation set one night later, a moth that stayed on the sheet, and removing them changes no score by more than 0.0003.
% TODO: re-check the 0.0003 in the mini_metrics convention once the rescoring lands.
